\section{Methods}
\label{sec:methods}

Notation: participant $p \in \{101,\dots,110\}$; crossing $i$ of participant $p(i)$ with feature vector
$\mathbf{x}_i \in \mathbb{R}^d$ and label $y_i$ (height H1/H2/H3, width W1/W2, or one of the six obstacles);
$O_p$ = obstacle crossings of $p$, $L_p$ = level walks of $p$; $\ell_p$ = leg length.

\subsection{Evaluation protocol}
\label{sec:eval}

\paragraph{Leave-one-participant-out (LOPO).} For each participant $p$, every model is trained on the crossings of the nine
other participants and predicts the crossings of $p$. All fitted steps (median imputation of missing values, standardisation,
feature selection, the model itself, thresholds) use the training participants only, so the reported accuracy is that of
a \emph{new person}. The ten sets of predictions are pooled before the metrics are computed
(\vn{mlTools.lopo}).

\paragraph{Balanced accuracy.} With $K$ classes and $n_k$ crossings of class $k$,
\begin{equation}
\mathrm{BA} = \frac{1}{K}\sum_{k=1}^{K} \frac{\#\{i : y_i = k,\ \hat y_i = k\}}{n_k},
\end{equation}
the mean per-class recall; chance is $1/K$ whatever the class balance.

\paragraph{Confidence intervals.} Participant bootstrap: the ten participants are drawn with replacement 2000 times, the
balanced accuracy of the pooled predictions of the drawn participants is recomputed, and the 2.5th and 97.5th percentiles
form the 95\,\% CI. A difference between two models uses the same draws for both (paired difference)
(\vn{mlTools.bootCI}).

\paragraph{Permutation tests.} The labels are shuffled \emph{within each participant} (keeping each person's class counts),
the complete LOPO pipeline is refitted, and this is repeated $n$ times; $p = (1 + \#\{\mathrm{BA}_{\text{null}} \geq
\mathrm{BA}_{\text{obs}}\}) / (1 + n)$ (\vn{mlTools.permTest}).

\paragraph{Nested (honest) selection.} When several settings are compared (feature set $\times$ normalisation $\times$
model $\times$ limb fusion), picking the best on the test results is optimistic. In the nested version the setting is chosen
inside each outer fold by an inner LOPO on the nine training participants, and only the chosen setting is applied to the
left-out participant.

\subsection{Normalisation}
\label{sec:norm}

\paragraph{Generic (level-walk) scaling.} Each feature is referenced to the person's normal walking, and lengths are
divided by the leg length:
\begin{equation}
\tilde x_{ij} = \frac{x_{ij} - \operatorname{median}_{l \in L_p} x_{lj}}{s_j}, \qquad s_j = \begin{cases} \ell_p\ (\text{mm}) & \text{if feature } j \text{ is a length}\\ 1 & \text{otherwise.}\end{cases}
\end{equation}
A device would get this from a short normal walk (\vn{mlTools.calibrate}).

\paragraph{Person scaling.} Each feature is scaled with the person's own obstacle crossings, without their labels:
\begin{equation}
z_{ij} = \frac{x_{ij} - \operatorname{median}_{k \in O_p} x_{kj}}{\operatorname{IQR}_{k \in O_p} x_{kj}},
\end{equation}
with the inter-quartile range from linear-interpolation quantiles (set to 1 if zero) (\vn{mlTools.selfnorm}). This removes
the person's walking ``signature'' and personal obstacle margin. It assumes the person's crossings cover several heights.

\paragraph{Causal person scaling (real time).} Crossing number $j$ of a session is scaled with the median and IQR of the
crossings \emph{before} it; until $n_{\min}$ crossings have been seen (5 or 10), the generic model is used.

\paragraph{Person-centring (haptic thresholds).} For single-signal rules, a value is expressed in physical units either
relative to the person's normal walking ($x - \operatorname{median}_{L_p} x$) or relative to the person's typical crossing
($x - \operatorname{median}_{O_p} x$), at the same moment of the swing or the same time before the trailing lift-off.

\subsection{Feature construction}
\label{sec:features}

\paragraph{Clearance-shape features} (\vn{mlTools.shapeFeats}). From one limb's foot-height curve $h(t)$ and progression
$x(t)$ over the swing (51 points): peak height $h^*$ and the progression to it; for $q \in \{0.5,0.7,0.8,0.9\}$ the horizontal
distance and the \% of swing during which $h \geq q\,h^*$; the foot tilt at the apex; the progression after the apex; and the
ratio of the steepest rise to the steepest fall of $h$. They describe how long and how far the foot stays high.

\paragraph{Time-resolved features} (\vn{mlTools.cutFeatures}). For data up to $c\,\%$ of the leading swing: for each of the six
signals, the value at $c$ and the maximum, minimum, mean and final slope from toe-off to $c$, plus the approach strides.

\paragraph{Stride and online features} (\vn{mlTools.strideTable}). For every ZVP-to-ZVP stride of both feet: peak height and
its timing, time and distance above 50/70/90\,\% of the peak, stride length, peak vertical speed, end height and duration;
and six features known at the foot apex (peak height, time from the stride start, distance travelled, maximum and mean
vertical speed before the apex, height 0.1\,s before the apex).

\paragraph{Out-of-plane features} (\vn{crossing_ml_lateral.m}). Frontal tilt (Euler $Y$) and heading (unwrapped Euler $Z$) of
each segment in the swing (range, largest deviation from toe-off, signed mean deviation, sign flipped for the right leg), and
the foot's sideways path from the foot trajectory.

\subsection{Models}
\label{sec:models}

\paragraph{L2 logistic regression} (\vn{mlTools.fitPipe}, kind \vn{logreg}). Inputs median-imputed and z-scored with the
training statistics, optionally reduced to the $k$ features with the largest ANOVA $F$. Weights $W$, intercepts $b$:
\begin{equation}
\min_{W,b}\ C \sum_i w_i\, \mathcal{L}\big(y_i, \operatorname{softmax}(W^\top \mathbf{z}_i + b)\big) + \tfrac12 \lVert W\rVert^2,
\qquad C = 0.1,
\end{equation}
with the cross-entropy loss $\mathcal{L}$ (sigmoid for two classes), solved with a quasi-Newton optimiser; intercepts are not
penalised (the scikit-learn objective, so the MATLAB and Python versions agree).

\paragraph{Random forest} (kind \vn{forest}). MATLAB \vn{TreeBagger}: 300--500 classification trees, minimum leaf size 2,
$\lfloor\sqrt d\rfloor$ candidate features per split, uniform class prior (balanced classes), fixed seed. Class probability =
share of trees voting for the class.

\paragraph{Support vector machine and boosting} (kinds \vn{svm}, \vn{boost}). RBF-kernel SVM (one-versus-one ECOC, automatic
kernel scale, box constraint 1) on z-scored inputs; AdaBoostM2 (three classes) or GentleBoost (two classes) with 200 trees
of at most 8 splits and learning rate 0.1. Their scores are turned into class probabilities with a softmax.

\paragraph{Separate limbs and soft vote.} The leading and trailing limbs are modelled separately (their mechanics differ);
the crossing is classified from the average of the two models' class probabilities,
$P(y \mid \text{crossing}) = \tfrac12\,[P_{\text{lead}}(y) + P_{\text{trail}}(y)]$. Ensembles average several models the same
way.

\paragraph{1D convolutional network} (\vn{crossing_ml_story.m}, part \vn{cnn}). Input: one limb's six swing curves
($6 \times 51$), normalised per participant and channel. Architecture: conv1d($5 \times 16$) -- batch normalisation -- ReLU --
conv1d($5 \times 32$) -- batch normalisation -- ReLU -- global average pooling -- dropout 0.3 -- fully connected (3) --
softmax. Training: Adam, 60 epochs, mini-batch 32, learning rate $10^{-3}$, L2 $10^{-3}$, training set tripled with Gaussian
noise ($\sigma = 0.1$); three seeds averaged; leading and trailing networks fused by soft vote.

\paragraph{Regression of the height in centimetres.} Target = nominal height $h_i = (\text{level}\,\% / 100)\,\ell_p$.
Ridge regression on z-scored features, $\hat\beta = (A^\top A + \lambda R)^{-1} A^\top h$ with $\lambda = 10$ (intercept not
penalised), or a random-forest regression (300 trees, minimum leaf 3). The predicted centimetres are turned into a class by
the nearest of the person's three nominal heights.

\paragraph{Physics-informed offset model} (\vn{crossing_ml_device.m}, D2). Physics: the foot lift rises by about the obstacle
height, but each person has a personal offset. A height score $s = \mathbf{w}^\top\mathbf{x}$ is learned by ridge regression
($\lambda = 10$) of the person-centred height level $t$ (\% leg length) on person-centred features of the training
participants, rescaled so that one score unit equals one \% of leg length. For a new person,
\begin{equation}
s_i = c + t_i + \varepsilon_i,\qquad \varepsilon_i \sim \mathcal{N}(0,\sigma^2),\qquad t_i \in \{10, 20, 30\},
\end{equation}
with one unknown offset $c$. The offset is set from a prior (population mean, or the level-walk score plus a learned shift,
prior SD $\tau$), or from $n$ labelled crossings as the posterior mean
\begin{equation}
\hat c = \frac{\sum_{i=1}^{n} (s_i - t_i)/\sigma^2 + \mu/\tau^2}{n/\sigma^2 + 1/\tau^2},
\end{equation}
or from unlabelled crossings with a constrained Gaussian mixture (components at $c + 10$, $c + 20$, $c + 30$, common $\sigma$,
free class shares, EM for 100 iterations). A crossing is classified by the nearest component, summing the log-likelihoods
of the leading- and trailing-foot scores.

\paragraph{Single-signal threshold rules (haptic).} Two ordered thresholds $t_1 \leq t_2$ on one signal $v$ at one moment
(direction $s = \operatorname{sign}\operatorname{corr}(v, \text{level})$): H1 if $s v < t_1$, H2 if $t_1 \leq s v < t_2$, H3
otherwise; $t_1, t_2$ maximise the training balanced accuracy over a grid of percentiles. For crossing detection, one
threshold maximising the Youden index (sensitivity + specificity $- 1$). The moment is chosen inside the training
participants (nested), optionally only among moments available before a deadline.

\subsection{Statistics and descriptive measures}
\label{sec:stats}

\begin{itemize}
  \item \textbf{Mixed-effects scaling.} $y = \beta_0 + \beta_1 h + u_{0p} + u_{1p} h + \varepsilon$ (MATLAB \vn{fitlme},
  random intercept and slope per participant, Satterthwaite degrees of freedom); $\beta_1$ = change per cm of obstacle.
  \item \textbf{Within-person effect size.} $\eta^2 = SS_{\text{between levels}} / SS_{\text{total}}$ after removing each
  participant's mean, computed at every 2\,\% of the swing (map) or per feature.
  \item \textbf{Separability.} $d' = \overline{\Delta\mu}/\sigma_w$: mean step between the mean foot lifts of neighbouring
  heights divided by the pooled within-height SD, per participant.
  \item \textbf{Map of all crossings.} t-SNE (perplexity 30) of the 110 features; silhouette score by participant and by
  height.
  \item \textbf{Calibration of probabilities.} Brier score $\frac1n\sum_i\sum_k (P_{ik} - [y_i = k])^2$ and expected calibration
  error $\sum_b \frac{n_b}{n}\,\lvert \mathrm{acc}_b - \mathrm{conf}_b\rvert$ over ten bins of the top-class probability.
  \item \textbf{Permutation importance.} In each left-out participant, one feature (or feature group) is shuffled across
  that participant's crossings; the drop in balanced accuracy, averaged over 10 shuffles, measures how much the model
  needs it.
\end{itemize}

\subsection{Analyses}
\label{sec:analyses}

Table~\ref{tab:analysis-map} maps the research questions to the nine MATLAB scripts in \vn{Analysis/ML/}; all share the
building blocks of \vn{mlTools.m}. Their reports (\vn{ml_*_matlab.md}) and figures (\vn{fig_*.png}) are in
\vn{Results/ML Models/MATLAB/}.

\begin{table}[htbp]
\centering
\caption{Analyses, the questions they answer and how.}
\label{tab:analysis-map}
\footnotesize
\begin{tabular}{@{}p{3.1cm} p{0.7cm} p{11.3cm}@{}}
\toprule
\textbf{Script} & \textbf{Q} & \textbf{What it does} \\
\midrule
\vn{crossing_ml_matlab.m} & Q2, Q3 & Height, width and the six obstacles from all / leading / trailing / segment features,
raw vs level scaling, logistic vs forest, top-$k$ features, permutation tests; cueing window and time-resolved prediction;
separate limb models and soft vote; personal calibration; nested selection; comparison with the earlier Python version. \\
\vn{crossing_ml_explore.m} & Q2 & Clearance-shape features and full curves; SVM, boosting and ensembles; ranking within a
person; self-training; personal calibration with 1--3 labelled crossings; width tricks; nested selection over all settings. \\
\vn{crossing_ml_story.m} & Q1, Q2 & $\eta^2$ maps over the swing; the personal signature (person identification); dose-response
and mixed-effects slopes; personal margin; height in cm; t-SNE; 1D CNN; width in the camera foot placement; clearance;
lead-limb choice and limb coupling. \\
\vn{crossing_ml_device.m} & Q3 & Which sensors; how much calibration (offset model, labelled and unlabelled crossings); answering only
when confident; learning curve over the number of training participants. \\
\vn{crossing_ml_stream.m} & Q3 & Camera-free pipeline on all strides: detect crossing strides, pair them into lead and trail,
classify the height; scored against the camera. \\
\vn{crossing_ml_realworld.m} & Q3, Q4 & Mixed obstacle order with causal person scaling (30 shuffles); a virtual real-time cue
at the leading-foot apex or landing; why some participants are harder ($d'$). \\
\vn{crossing_ml_lateral.m} & Q2 & Whether width shows in the side-to-side (out-of-plane) motion. \\
\vn{crossing_ml_validate.m} & Q1--Q3 & Unseen heights and widths; what sets the clearance; permutation importance; paired
statistics, confusion matrix and probability calibration; anticipation in the approach; IMU vs camera scaling. \\
\vn{crossing_ml_haptic.m} & Q4 & Event timeline relative to the trailing toe lift-off; single-signal crossing detection and
height thresholds over the swing and in real time; extra angle and lift per height. \\
\bottomrule
\end{tabular}
\end{table}

\paragraph{Generalisation tests} (\vn{crossing_ml_validate.m}). \emph{Unseen height}: within LOPO, the regression is trained on
the other participants' crossings of two height levels only and tested on the left-out participant's crossings of the third
level (H2 left out = interpolation; H1 or H3 left out = extrapolation). \emph{Unseen width}: height recognition trained on the
other participants' crossings of one width only and tested on the other width, compared with training on the same width
(same amount of training data).

\paragraph{Clearance mechanism} (\vn{crossing_ml_validate.m}, V3). Within each participant $\times$ obstacle cell (cell means
removed), the share of the crossing-to-crossing clearance variance explained by the camera toe peak above the obstacle, by the
camera take-off distance (toe before the obstacle at its toe-off), by the IMU foot lift and by the IMU toe estimate, and the
absolute clearance of a new participant from ordinary least squares (LOPO).

\paragraph{Real-time cue} (\vn{crossing_ml_realworld.m}, R2). Crossing strides are detected online at the leading-foot apex
from the six apex features (all strides of the session); the height is predicted from the leading-limb data available at
the leading-foot apex (50\,\% of its swing) or at its landing (100\,\%), with causal person scaling; the cue is given always or
only if the top probability is at least 0.6. Outcome per crossing: correct cue, wrong cue or no cue.

\paragraph{Haptic timing and thresholds} (\vn{crossing_ml_haptic.m}). Event times relative to the camera trailing toe lift-off;
single-signal detection (crossing vs level stride) and three-level grading over the swing; the same grading on the
time-aligned signals for cue deadlines of 500, 300, 200, 100 and 0\,ms before the lift-off; the mean extra angle or lift per
height at the most informative moment, and its crossing-to-crossing SD.
